\documentclass[11pt]{article}
\usepackage[margin=1in]{geometry}
\usepackage{amsmath,amssymb,amsthm}
\usepackage{hyperref}
\newtheorem{theorem}{Theorem}
\title{M11: Two-user NOMA with arbitrary fixed QoS}
\author{Internal benchmark model; selected idea-01}
\date{P4 draft for independent review and G4; no simulation results}
\begin{document}
\maketitle

\section{Scope, normalization and physical model}
This is an established-method dry run under approved G1--G3, not an
originality claim. Every statement below concerns a single constant-channel
block. Parameters are finite real numbers with
\(P\geq0\), \(0\leq g_w\leq g_s\), and fixed individual targets
\(q_w,q_s\geq0\) in bit/s/Hz over the whole block. No fading distribution,
physical deployment, numerical tolerance or experimental protocol is selected.
Noise variance and bandwidth are normalized to one. Thus \(P\) is linear
transmit power relative to normalized noise, not a dBm input or a physical
power value in watts; \(A=Pg_w\), \(B=Pg_s\) are dimensionless received
SNRs if a user receives full power. No unreported path-loss or SNR conversion
is used. Let scalar complex channels satisfy \(g_i=|h_i|^2\), for
\(i\in\{w,s\}\). Perfect instantaneous CSI and ideal SIC are scoped
idealizations. Independent unit-energy Gaussian codebook symbols \(s_w,s_s\)
give the information-theoretic rate model, not a finite-modulation link model.
Receiver noise \(z_i\sim\mathcal{CN}(0,1)\) has real and imaginary
variance \(1/2\). The full-range strong-user power fraction is
\(a\in[0,1]\); its weak-user counterpart is \(1-a\). The signal equations are
\begin{align}
 x&=\sqrt{P(1-a)}s_w+\sqrt{Pa}s_s, & y_i&=h_ix+z_i.
\end{align}
The average transmitted power is \(P\). An endpoint can carry a zero-rate
message with zero allocated power. There is no constraint \(a\leq1/2\),
SIC margin, residual interference, adaptive time split or altered objective.

\section{Both weak-message decoding sites and rates}
The weak receiver treats the strong message as noise. The strong receiver
first decodes the weak message while treating its own message as noise,
then removes the decoded weak message exactly. Their weak-message SINRs are
\begin{align}
 \gamma_{w\to w}(a)&=\frac{(1-a)A}{1+aA}, &
 \gamma_{w\to s}(a)&=\frac{(1-a)B}{1+aB}, &
 \gamma_{s\to s}(a)&=aB.
\end{align}
In particular,
\begin{equation}
 \gamma_{w\to s}-\gamma_{w\to w}
 =\frac{(1-a)(B-A)}{(1+aA)(1+aB)}\geq0.
\end{equation}
The weak code rate must not exceed either weak-message decoding capacity;
their minimum equals the weak receiver's capacity even at equal gains and
endpoints. Hence the achievable boundary rates and objective are
\begin{align}
 R_w(a)&=\log_2\frac{1+A}{1+aA}, &R_s(a)&=\log_2(1+aB),\\
 F(a)&=R_w(a)+R_s(a).
\end{align}
Choosing a smaller actual code rate is permitted, but cannot improve the
sum-rate maximum. Capacity equalities are ideal asymptotic achievable-rate
limits; no finite-block error probability or practical SIC claim is inferred.
Both target inequalities use \(\geq\), so equality is feasible in exact
real arithmetic.

\section{Exact QoS feasible set, including zero cases}
Define \(T_i=2^{q_i}-1\). For \(B>0\), strong QoS is exactly
\(a\geq T_s/B\). For \(A>0\), weak QoS is exactly
\begin{equation}
 (1-a)A\geq T_w(1+aA)
 \quad\Longleftrightarrow\quad
 a\leq\frac{A-T_w}{A(1+T_w)}.
\end{equation}
These follow by applying a strictly increasing logarithm and multiplying
strictly positive denominators; SIC imposes no additional bound because
the second weak decoding site has no smaller capacity.
Use the following branches, without division by zero:
\begin{itemize}
 \item If \(B=0\), ordering also gives \(A=0\). Feasible iff
 \(q_w=q_s=0\), with feasible set \([0,1]\); otherwise infeasible.
 \item If \(A=0<B\), feasible weak QoS requires \(q_w=0\).
 Then set \(L=T_s/B\), \(U=1\).
 \item If \(A>0\) (therefore \(B>0\)), set
 \(L=T_s/B\), \(U=\min\{1,(A-T_w)/(A(1+T_w))\}\).
\end{itemize}
For either positive-\(B\) branch, the set is \(\mathcal A=[L,U]\)
iff \(L\leq U\); otherwise it is empty. Nonnegative targets ensure
\(L\geq0\) and \(U\leq1\). A negative \(U\), \(L>1\), or
\(L>U\) is infeasibility, not an instruction to clamp an answer into
\([0,1]\). For \(A>0\) an equivalent complete feasibility condition is
\begin{equation}
 \frac{T_s}{B}+\frac{T_w}{A}+\frac{T_wT_s}{B}\leq1.
\end{equation}
A singleton \(L=U\) is feasible. Zero total power, zero strong gain,
and both zero gains all enter the \(A=B=0\) branch. Zero weak gain at
positive received strong SNR enters the second branch. Zero QoS gives
the corresponding inactive bound: \(T_s=0\) implies \(L=0\);
\(T_w=0\) implies \(U=1\). Arbitrarily large finite targets can simply
make the interval empty; they are not restricted to paper-specific regimes.

\section{Optimizer, curvature and degeneracies}
\begin{theorem}[Exact scalar solution]
If the feasible set is nonempty, \(a^\star=U\) is a maximizer (use
\(U=1\) in the all-zero branch). If \(B>A\), it is the unique maximizer.
If \(B=A\), every feasible split is optimal. An empty set has no optimizer
and no constrained optimum sum rate.
\end{theorem}
\begin{proof}
On \([0,1]\), direct differentiation gives
\begin{align}
 F'(a)&=\frac{1}{\ln2}\left(\frac{B}{1+aB}-\frac{A}{1+aA}\right)
       =\frac{B-A}{\ln2(1+aA)(1+aB)},\\
 F''(a)&=\frac{1}{\ln2}\left(\frac{A^2}{(1+aA)^2}
                              -\frac{B^2}{(1+aB)^2}\right)\\
 &=-\frac{(B-A)(A+B+2aAB)}{\ln2(1+aA)^2(1+aB)^2}.
\end{align}
Positive denominators give monotonicity; if \(B>A\), then \(B>0\)
and both the first derivative is strictly positive and second derivative
strictly negative. Consequently the largest feasible split is the unique
maximum, including a singleton. If \(B=A\), \(F(a)=\log_2(1+A)\)
is constant. This includes \(A=B=0\), when every rate is zero.
\end{proof}
For \(A>0\), the selected endpoint exactly meets weak QoS,
\(R_w(U)=q_w\), including \(q_w=0,U=1\). For \(A=0<B\),
\(U=1\) gives weak rate zero and strong rate \(\log_2(1+B)\).
Equal gains do not imply equal user rates or equal targets; they only
make the sum constant on the feasible interval. For \(A=B>0\), feasibility
is equivalently \(q_w+q_s\leq\log_2(1+A)\).

The reduced optimization is maximization of a concave scalar function over
an interval, equivalently convex minimization of \(-F\). Strict concavity
holds only for \(B>A\). The weak rate alone is convex and nonincreasing
in \(a\), and the strong rate is concave and nondecreasing; their QoS
sets become affine interval inequalities by exact algebra. No claim of
joint convexity in channels, targets or total power is made. There are
no alternating blocks or iterative convergence claims. Endpoint evaluation
requires a constant number of real-arithmetic operations and elementary
functions, hence \(O(1)\) work and storage per block in that model.
This is not constant bit complexity or constant floating-point latency.

\section{Matched equal-time OMA and conditional comparison}
OMA serves each user for half the same block, using all the same bandwidth
and transmit power \(P\) while active, with the same channels and targets.
It has
\begin{align}
 R_w^{O}&=\tfrac12\log_2(1+A), &R_s^{O}&=\tfrac12\log_2(1+B), &
 F^O&=R_w^O+R_s^O.
\end{align}
OMA is feasible iff \(R_i^O\geq q_i\) for both users, equivalently
\(A\geq2^{2q_w}-1\) and \(B\geq2^{2q_s}-1\).
The baseline rate is per whole block, not per active slot. If block duration
is \(\tau\), both schemes spend normalized energy \(P\tau\):
NOMA transmits at \(P\) throughout, OMA uses \(P\) in each disjoint
half-slot. Neither OMA user uses \(2P\) or \(P/2\). Both occupy one
bandwidth-block resource; noise bandwidth stays one during an active slot.
No optimized OMA scheduling or hardware/CSI/SIC overhead adjustment is included.

\begin{theorem}[Scope-wide OMA inclusion and sum comparison]
With this model and matched reference, OMA feasibility implies NOMA
feasibility. On the both-feasible set, \(F(a^\star)\geq F^O\).
Equality occurs iff \(A=B\); if \(B>A\), the sum inequality is strict.
These conclusions need no fading distribution or high-SNR approximation.
\end{theorem}
\begin{proof}
First suppose \(A>0\). Let \(c(v)=1/(1+\sqrt{1+v})\) for
\(v\geq0\). This function is decreasing, so \(c(B)\leq c(A)\).
Solving the two individual inequalities against their OMA rates gives
\(a\leq c(A)\) for the weak user and \(a\geq c(B)\) for the strong.
The witness \(a_F=c(A)\) therefore attains
\(R_w(a_F)=R_w^O\) and \(R_s(a_F)\geq R_s^O\).
If OMA meets the fixed targets, this same witness meets them too, so it
belongs to \(\mathcal A\). Optimality gives
\(F(a^\star)\geq F(a_F)\geq F^O\). When \(B>A\),
\(c(A)>c(B)\) and the strong inequality is strict since \(B>0\).
For \(A=0<B\), OMA feasibility forces \(q_w=0\), and the separate
witness \(a_F=1\) gives \(R_w=R_w^O=0\) and
\(R_s=\log_2(1+B)>R_s^O\); it also meets \(q_s\).
For \(A=B=0\), OMA feasibility requires both targets zero and the
all-zero NOMA branch is feasible with equal zero sums. Finally if
\(A=B>0\), every feasible NOMA sum is \(\log_2(1+A)=F^O\).
\end{proof}
The fairness witness is a proof device, not a substituted objective,
new QoS choice, or selected implementation split. A sum-optimal split may
give the weak user less than its OMA rate when its fixed target is lower
than that OMA rate. No individual-rate dominance is asserted for the
optimizer. The theorem says nothing about arbitrary fixed NOMA splits,
unmatched OMA, imperfect SIC/CSI, finite-block coding or other scopes.

\section{Feasibility ledger and prospective validation}
For a later ensemble of \(N\) matched blocks, let \(I_N(j),I_O(j)\)
be feasibility indicators, and \(n_{uv}\) count blocks with
\((I_N,I_O)=(u,v)\), \(u,v\in\{0,1\}\). Report all four counts,
\(N=\sum_{u,v}n_{uv}\), NOMA fraction \((n_{10}+n_{11})/N\),
OMA fraction \((n_{01}+n_{11})/N\), and paired count \(n_{11}\).
The exact theorem predicts \(n_{01}=0\); a numerical positive count
requires diagnosis, not deletion. Fractions for \(N=0\) are undefined.
Rate differences \(\Delta(j)=F(a^\star(j))-F^O(j)\) and their paired
summary use only \(n_{11}\), with that denominator explicitly reported;
if \(n_{11}=0\), the paired summary is undefined, not zero.
For infeasible NOMA, retain a failure status and undefined constrained
optimum; do not substitute zero rate. OMA's unconstrained reference rates
may still be recorded outside its feasible set, explicitly labeled, but
do not enter the accepted feasible-case comparison.

All identities and inequalities here are exact derivations. The bounded
symbolic file checks selected algebraic identities, limits and decidable
curvature cases; it cannot prove every inequality, domain branch, source
claim, or future implementation. Independent scalar-oracle validation is
proposed in algorithm.md, not executed. Numerical cancellation in
\(2^q-1\), subtraction near \(A=T_w\), overflow in products/exponentials,
and misclassification near a singleton interval must be addressed in P5.
Stable log1p/expm1 forms, scaled comparisons and higher precision are
prospective tools; tolerances, resolutions, caps and stopping criteria
remain unstated pending P5 approval. Never infer measured results from
this specification.

\section{Source attribution}
The approved scope and selected idea determine the model. Source-map.md
links the six grounded cached manuscripts with PDF page numbers and transfer
restrictions. Xu et al. (arXiv:1504.07751, pp.1--2) supports the ordered
Gaussian broadcast rate context and warns through fixed-split loss events.
Zhu et al. (DOI:10.1109/JSAC.2017.2725618, pp.3,6--7) supplies a related
boundary optimizer with extra power-order/SIC-stability restrictions that
are not imported. Oviedo and Sadjadpour's Fair-NOMA manuscript
(DOI:10.1109/TVT.2017.2689000, pp.2--3) supplies the matched OMA and
individual fairness interval context. The exact arbitrary-QoS derivation,
zero-case extensions and scoped comparison proof above are this project's
own rederivations of established principles, not copied project results
or an original contribution. No asymptotic numerical gain is asserted.
\end{document}
